\section{Approach}\label{section:approach}

\newcommand{\promptin}[1]{\textcolor{blue!70!black}{{#1}}}
\newcommand{\promptout}[1]{\textcolor{red!70!black}{{\ul{#1}}}}

\begin{figure}[!t]
\centering\small
\scalebox{0.83}{
\begin{tabular}{@{}p{1.18\columnwidth}@{}}
\toprule
\multicolumn{1}{c}{(a) \textbf{Query generation} $x \to \{q_1, \dots, q_N\}$} \\
You said: \promptin{Your nose switches back and forth between nostrils. When you sleep, you switch about every 45 minutes.  This is to prevent a buildup of mucus. It's called the nasal cycle.} \\
To verify it, \\
\promptout{a) I googled: Does your nose switch between nostrils?} \\
\promptout{b) I googled: How often does your nostrils switch?} \\
\promptout{c) I googled: Why does your nostril switch?} \\
\promptout{d) I googled: What is nasal cycle?} \\
\midrule
\multicolumn{1}{c}{(b) \textbf{Agreement model} $(y, q, e) \to \{0, 1\}$} \\
You said: \promptin{Your nose switches \dots(same as above)\dots nasal cycle.} \\
I checked: \promptin{How often do your nostrils switch?} \\
I found this article: \promptin{Although we don't usually notice it, during the nasal cycle one nostril becomes congested and thus contributes less to airflow, while the other becomes decongested. On average, the congestion pattern switches about every 2 hours, according to a small 2016 study published in the journal PLOS One.} \\
\promptout{Your nose's switching time is about every 2 hours, not 45 minutes.} \\
\promptout{This disagrees with what you said.} \\
\midrule
\multicolumn{1}{c}{(c) \textbf{Edit model} $(y, q, e) \to \text{new }y$} \\
You said: \promptin{Your nose switches \dots(same as above)\dots nasal cycle.} \\
I checked: \promptin{How often do your nostrils switch?} \\
I found this article: \promptin{Although we \dots(same as above)\dots PLOS One.} \\
\promptout{This suggests 45 minutes switch time in your statement is wrong.} \\
\promptout{My fix: Your nose switches back and forth between nostrils. When you sleep, you switch about every 2 hours. This is to prevent a buildup of mucus. It's called the nasal cycle.} \\
\bottomrule
\end{tabular}
}
\begin{subfigure}{0pt}\phantomsubcaption\label{fig:example_prompts_cqgen}\end{subfigure}
\begin{subfigure}{0pt}\phantomsubcaption\label{fig:example_prompts_agreement}\end{subfigure}
\begin{subfigure}{0pt}\phantomsubcaption\label{fig:example_prompts_revision}\end{subfigure}
\caption{
\label{fig:example_prompts}
\textbf{Examples of few-shot examples} used to prompt the PaLM model (\promptin{blue} = input; \promptout{red} = output).
}
\end{figure}

We now present \approachfull{} (\approach{}), a simple method for solving the \taskName{} task.
As illustrated in Figure~\ref{fig:overview}, given an input passage $x$,
the research stage first generates a set of queries $\{q_1, ..., q_N\}$, each investigating one aspect of $x$ that potentially requires attribution.
For each query $q_i$, it retrieves web documents and selects the best evidence snippets $\{\e_{i1}, \e_{i2},\dots\}$.
    The revision stage then revises the original text $x$ using the retrieval results $\{(q_1, \e_{11}), \dots\}$, yielding a revised text $y$.

Most components for \approach{} are implemented using few-shot prompting \cite{Brown2020LanguageMA}. We use PaLM \cite{Chowdhery2022PaLMSL} as our language model.
Figure~\ref{fig:example_prompts} shows some few-shot examples we use, while Appendix~\ref{section:prompting} lists the full prompts.





\subsection{Research stage}
\label{section:qgen}


\paragraph{Query generation}
We perform \emph{comprehensive question generation} (CQGen) which produces a sequence of questions covering \emph{all aspects} of the passage $x$ that need to be verified and attributed.
A similar strategy has been employed to train text-planning models \cite{narayan2022blueprint}.
A prompt with six human demonstrations %
was sufficient for PaLM to adequately learn the task.  To increase diversity and coverage, we sample from our CQGen model three times and take the union of the resulting queries.


\paragraph{Evidence retrieval}
For each query from CQGen, we use Google Search to retrieve $K=5$ web pages.
We extract candidate evidence snippets from each web page by running a sliding window of four sentences across the page, breaking at document headings.
The evidence snippets for each query are then ranked based on their relevance to the query. For this, we use an existing query-document relevance model trained following \citet{ni2021gtr}, which computes a relevance score $S_{\text{relevance}}(q, e)$ between a query $q$ and an evidence snippet $e$. We then keep the top $J=1$ evidence for each query. 
The final retrieval result is
$[
(q_1, \e_{11}),\allowbreak
\ldots,\allowbreak
(q_1, \e_{1J}),\allowbreak
\ldots,\allowbreak
(q_N, \e_{N1}),\allowbreak
\dots,\allowbreak
(q_N, \e_{NJ})]$,
where $\e_{ij}$ denotes the $j^{\text{th}}$ evidence for the $i^{\text{th}}$ query, and $N$ denotes the total number of queries from CQGen (which can be different for each input $x$).


\subsection{Revision stage}\label{sec:approach-editor}

After retrieving evidence, certain parts of $x$ may now be properly attributed, but other parts remain unattributed and should be revised.
As illustrated in Figure~\ref{fig:overview},
the revision stage initializes the output $y = x$. Then for each retrieved $(q, e) = (q_i, \e_{ij})$,
the \emph{agreement model} checks if the evidence $e$ disagrees with the current output $y$ regarding the issue in query $q$. If a disagreement is detected, the \emph{edit model} edits $y$ to agree with $e$; otherwise, it does nothing. The process continues until all retrievals are processed.

\paragraph{Agreement model}
The agreement model takes the partially edited passage $y$, a query $q$, and the evidence $e$ as input. It then decides whether both $y$ and $e$ imply the same answer to the question in $q$. This form of question-guided agreement was previously explored by \citet{honovich2021q2}.
We implement this by few-shot prompting PaLM using a chain-of-thought style prompt \cite{Wei2022COT}, where we ask the model to explicitly state the implied answers for both $y$ and $e$ before producing its judgment about their agreement. %

\paragraph{Edit model}
The edit model is run only if a disagreement is detected.
The model takes $y$, $q$ and $e$ as input, and outputs a new version of $y$ that aims to agree with $e$ while otherwise minimally altering $y$. We again use few-shot prompting and chain-of-thought, where we ask the model to first identify a particular span in $y$ that needs to be edited before generating the revised $y$. %
This helps reduce the editor's deviation from the current $y$.\footnote{The editor occasionally produces large edits that bring the new revision close to $e$ but far from the current $y$. Since this is rarely desirable, we reject edits with edit distance above 50 characters or 0.5 times the original text length.}


\subsection{Attribution report}\label{sec:attribution-report}

Finally, we select at most $M=5$ evidence snippets to form an attribution report $A$. Note that during evidence retrieval and revision, we may have encountered and used more than $M$ snippets.
Our goal is to find a subset of snippets that maximizes \emph{coverage} over the potentially attributable points in the passage, as represented by the queries ${q_1,\dots,q_N}$. We use the relevance model from Section~\ref{section:qgen} as a proxy for measuring how much an evidence $e$ covers the point raised by a query $q$. Then, we exhaustively search for $A \subseteq \{\e_{11}, \dots, \e_{NJ}\}$ of size at most $M$ that maximizes
\begin{equation}
\text{Cover}(A, q_{1:N}) := \sum_{i=1}^N \max_{e \in A} S_{\text{relevance}}(q_i, e).
\end{equation}
